\documentclass{amsproc}
\usepackage{amssymb}
\usepackage[hidelinks]{hyperref}

\newtheorem{theorem}{Theorem}[section]
\newtheorem{lemma}[theorem]{Lemma}
\theoremstyle{definition}
\newtheorem{definition}[theorem]{Definition}
\theoremstyle{remark}
\newtheorem{remark}[theorem]{Remark}

\numberwithin{equation}{section}

\begin{document}

\title{Title of the Paper}

\author{First Author}
\address{Department of Mathematics, University, City, Country}
\email{first@example.org}

\author{Second Author}
\address{Institute, City, Country}
\email{second@example.org}

\subjclass[2020]{Primary 11A41; Secondary 11N05}
\keywords{first keyword, second keyword}

\begin{abstract}
State the main results of the paper in a few sentences.
\end{abstract}

\maketitle

\section{Introduction}
Introduce the problem and the main result~\cite{hardy2008introduction}.

\section{Main result}
\begin{definition}
An integer $p > 1$ is \emph{prime} if its only positive divisors are $1$ and $p$.
\end{definition}

\begin{theorem}\label{thm:main}
There are infinitely many primes.
\end{theorem}

\begin{proof}
Suppose $p_1, \dots, p_n$ were all of them, and let $N = p_1 \cdots p_n + 1$.
No $p_i$ divides $N$, so a prime factor of $N$ is missing from the list.
\end{proof}

\begin{remark}
Theorem~\ref{thm:main} also follows from the divergence of
\begin{equation}\label{eq:sum}
  \sum_{p} \frac{1}{p} = \infty .
\end{equation}
\end{remark}

\bibliographystyle{amsplain}
\bibliography{references}

\end{document}
